\begin{frame}{자주 하는 실수 5가지 — 16.2 리뷰의 ``반박 가능한 질문''이 된다}
  \small
  \begin{enumerate}
    \item \textbf{시드 1개 실험으로 결론} — 같은 PPO라도 시드만 바꿔도
      마지막 10 평균이 $\pm100\sim200$ 흔들림(실습표: $-668\sim-732$).
      알고리즘 간 차이가 시드 간 sd보다 작으면 결론은 ``유의미한 차이
      관측 불가''. \emph{리뷰 질문: ``시드 간 sd가 방법 간 차이보다 커요?''}
    \item \textbf{학습 도중 리턴을 ``최종 정책 성과''로 보고} — 학습 도중
      마지막 10($-673\sim-732$)은 $\sigma$ 샘플링 노이즈가 섞인 ``탐험
      포함 성과''. ``곡선이 올랐다''를 ``잘 학습됐다''로 번역하는 순간
      읽기 3의 격차(결정적 $-1403$ vs 베이스라인 $-1274$)가 사라진다.
    \item \textbf{결정적 평가가 예뻐질 때까지 하이퍼파라미터 조정} —
      학습 곡선으로 조정하는 것은 정당한 val 선택이지만, \textbf{최종
      평가치를 보며 조정하면 평가에 과적합}(ML1 6.2 ``test는 한 번만'').
      구분선 = 선택 기준의 \textbf{사전 기록}.
    \item \textbf{베이스라인 없이 ``잘 배웠다''} — 무작위 리턴($-1274$)을
      재지 않고 ``$-1400$, 꽤 좋다''라고 발표 = 그 숫자가 무작위보다
      \emph{나쁜}지도 모름. zero-torque($\approx-1230$)까지 함께 재면
      기준선 두 개가 생긴다.
    \item \textbf{``결정적 $\ll$ 샘플 평가는 버그''로 오해} — $a=\mu$가
      $a\sim\mathcal{N}$보다 낮은 것은 \textbf{정상}. 격차의 크기는
      ``정책이 탐험 운에 얼마나 의존하는가''의 측정값 — \textbf{작을수록}
      정책이 자기 $\mu$만으로 선다.
  \end{enumerate}
\end{frame}
